\section{Билет 2. Симплекс-метод. Искусственный базис}

\ticketq{Симплекс-метод. Метод искусственного базиса.}

\paragraph{Обозначения.}
Каноническая задача $c\T[N]x[N] \to \min$, $Ax = b$, $x \ge 0$. Метод строит
опорные векторы $x^0, x^1, \dots$ с невозрастающим $c\T x^k$ (оптимум
достигается в опорном векторе — билет~1).

На $k$-й итерации: базис $N_k$, $\abs{N_k} = m$, $\det A[M,N_k] \ne 0$;
$B[N_k,M] = A[M,N_k]^{-1}$; $L_k = N \setminus N_k$, $x^k[L_k] = 0$.
Потенциалы и оценки:
\[
  A\T[M,N_k]\,y^k[M] = c[N_k] ,
  \qquad
  d[j] = c[j] - (y^k)\T a_j[M] , \quad j \in L_k .
\]

\paragraph{Теорема (критерий оптимальности).}
\textit{$d[j] \ge 0$ $\forall j \in L_k$ $\Rightarrow$ $x^k$ оптимален.}

\textsc{Док-во.} $x \in S$. Из $Ax = Ax^k = b$, разбивая на базисную и
небазисную части:
\begin{gather*}
  A[M,N_k]\big(x[N_k] - x^k[N_k]\big) + A[M,L_k]x[L_k] = 0 \\
  \Rightarrow\quad x[N_k] - x^k[N_k] = -B[N_k,M]A[M,L_k]x[L_k] .
\end{gather*}
Подставляя и учитывая $(y^k)\T = c\T[N_k]B[N_k,M]$:
\[
  c\T x - c\T x^k
   = c\T[N_k]\big(x[N_k] - x^k[N_k]\big) + c\T[L_k]x[L_k]
   = \sum_{j \in L_k} d[j]\,x[j] \ \ge\ 0 . \qquad \blacksquare
\]

\paragraph{Направление перехода.}
Пусть $\exists\, j_k : d[j_k] < 0$. Берём
\[
  u^k[N_k] = B[N_k,M]a_{j_k}[M] , \quad u^k[j_k] = -1 , \quad
  u^k[L_k \setminus \{j_k\}] = 0 , \quad x_\theta = x^k - \theta u^k .
\]
\textbf{Равенства:} $Au^k = A[M,N_k]B[N_k,M]a_{j_k} - a_{j_k} = 0$ (так как
$A[M,N_k]B = E$) $\Rightarrow$ $Ax_\theta = b$ $\forall\theta$.

\textbf{Цель:} $c\T u^k = (y^k)\T a_{j_k} - c[j_k] = -d[j_k]$, откуда
\[
  F(x_\theta) = F(x^k) + \theta\,d[j_k] < F(x^k) \quad (\theta > 0) .
\]

\paragraph{Выбор шага.}
Требуется $x^k[i] - \theta u^k[i] \ge 0$, $i \in N_k$:
\[
  u^k[i] \le 0 \ \forall i \in N_k \ \Rightarrow\ F(x_\theta) \to -\infty
  \ \ (\text{не ограничена снизу});
\]
иначе $\theta_k = \min\limits_{i \in N_k,\ u^k[i]>0} x^k[i]/u^k[i]$ на
индексе $i_k$; тогда $x_{\theta_k}[i_k] = 0$: $x[i_k]$ выходит, $x[j_k]$ входит.

\paragraph{Алгоритм.}
Дан $x^0$ и базис $N_0$.
\begin{csteps}
  \item $B[N_k,M] = A[M,N_k]^{-1}$; найти $y^k$ из $A\T[M,N_k]y^k = c[N_k]$.
  \item $d[j] = c[j] - (y^k)\T a_j$, $j \in L_k$.
  \item Все $d[j] \ge 0$ — $x^k$ оптимален, конец.
  \item Иначе выбрать $j_k$ с $d[j_k] < 0$, построить $u^k$.
  \item $u^k[i] \le 0$ $\forall i \in N_k$ — не ограничена снизу, конец.
  \item Найти $\theta_k$, $i_k$; $x^{k+1} = x^k - \theta_k u^k$,
        $N_{k+1} = (N_k \setminus \{i_k\}) \cup \{j_k\}$, к~шагу~1.
\end{csteps}

\paragraph{Конечность.}
В невырожденном случае $\theta_k > 0$, и по формуле изменения цели
$F$ \emph{строго} убывает $\Rightarrow$ базисы не повторяются; базисов не
больше $\binom{n}{m}$ (лемма билета~1) $\Rightarrow$ алгоритм конечен.

\paragraph{Вырожденность и зацикливание.}
$\abs{\Nzp} < m$ $\Rightarrow$ возможно $\theta_k = 0$: вектор не меняется,
меняется базис, $F$ не убывает $\Rightarrow$ \textbf{зацикливание} (возврат к
прежнему базису). Средство — \textbf{правило Бленда}: при нескольких
кандидатах на ввод и на вывод брать наименьший индекс; тогда ни один базис не
повторяется и конечность восстанавливается.

\paragraph{Метод искусственного базиса.}
Нужен, когда $x^0$ неизвестен. Добиваемся $b[M] \ge 0$ (строку с $b[i] < 0$
умножаем на $-1$) и решаем
\[
  \sum_{i \in M} y[i] \to \min , \quad
  A[M,N]x[N] + E[M,M]y[M] = b[M] , \quad x \ge 0,\ y \ge 0 ,
\]
старт $x[N] = 0$, $y[M] = b[M]$ — очевидный опорный вектор с базисом $E$.

\paragraph{Теорема (об искусственном базисе).}
\textit{Пусть $(x^*, y^*)$ — оптимум вспомогательной задачи. Тогда
$\sum_i y^*[i] > 0$ $\Rightarrow$ исходная задача несовместна;
$\sum_i y^*[i] = 0$ $\Rightarrow$ $x^*[N] \in S$.}

\textsc{Док-во.} $y \ge 0 \Rightarrow \sum_i y[i] \ge 0$: минимум неотрицателен.

\emph{Случай $>0$}, от противного: $\exists\,\hat{x} \in S$ $\Rightarrow$ пара
$(\hat{x}, 0)$ допустима со значением $0$ — противоречие.

\emph{Случай $=0$:} $y^*[i] \ge 0$ и $\sum_i y^*[i] = 0$ $\Rightarrow$
$y^*[M] = 0$ $\Rightarrow$ $A[M,N]x^*[N] = b[M]$, а $x^* \ge 0$ $\Rightarrow$
$x^* \in S$. $\blacksquare$

\paragraph{Алгоритм искусственного базиса.}
\begin{csteps}
  \item Привести к $b[M] \ge 0$, ввести $y[M] \ge 0$, составить вспомогательную
        задачу.
  \item Решить её симплекс-методом от базиса $x = 0$, $y = b$.
  \item $\sum_i y^*[i] > 0$ — несовместна, конец.
  \item Иначе $y^* = 0$, $x^*$ допустим.
  \item Искусственные переменные, оставшиеся в базисе с нулём, заменить
        столбцами $A[M,N]$ с сохранением линейной независимости; если замены
        нет — уравнение зависимо, вычеркнуть.
  \item Полученный базис — начальный для симплекс-метода на исходной задаче.
\end{csteps}

\seefull{гл.~1, \S7, \S9, \S10 (численные примеры там же)}
